
\section{General Matrix Solution of the Propagation-Orthogonality Constraint}\label{appx_general_matrix_source}
In the adapted basis, let $U=\rho(\mathbf{U})$, $\hat{U}=U/c$, and $R=\rho(\mathbf{R})$, with $\mathbf{U}^{2}=c^{2}$ and $\mathfrak{i}=\rho(\mathbf{i})$. The PO-source conditions are
\begin{equation}\label{eq_source_admissibility}
    \mathbf{U}\cdot\mathbf{R}=0,\quad R^{\ddagger}=R,\quad R'=gRg^{-1},
\end{equation}
where $g=\rho(\mathbf{g})$ represents an admissible local field-frame transformation and
\begin{equation}
    X^{\ddagger}\coloneqq\Gamma^{0}X^{\dagger}\Gamma^{0}, \quad \Gamma^{0}=\sigma_z\otimes\gamma^{0}.
\end{equation}
For a homogeneous grade-$r$ component $\mathbf{R}_r$, the contraction is given by
\begin{equation}
    \mathbf{U}\cdot\mathbf{R}_r=\frac{1}{2}\left(\mathbf{U}\mathbf{R}_r-(-1)^r\mathbf{R}_r\mathbf{U}\right),
\end{equation}
which shows that propagation orthogonality requires commutation for even grades and anticommutation for odd grades. Because contraction lowers grade, these conditions hold separately for each component of an inhomogeneous PO-source, with the scalar component satisfying the constraint identically.

In the chosen matrix realization, grade involution is represented by
\begin{subequations}
    \begin{align}
        &X\mapsto KXK, \quad K=\sigma_z\otimes\gamma_5,\\
        &\{K,\hat{U}\}=0, \quad [K,\mathfrak{i}]=0.
    \end{align}
\end{subequations}
The represented constraint becomes
\begin{equation}\label{eq_matrix_source_constraint}
    \rho\!\left(\mathbf{U}\cdot\mathbf{R}\right)=\frac{1}{2}(UR-KRKU)=0,
\end{equation}
whose general solution, for an arbitrary seed matrix $A\in M_8(\mathbb{C})$, is
\begin{equation}\label{eq_general_matrix_source}
    R=P(A)\coloneqq\frac{1}{2}\left(A+\hat{U}KAK\hat{U}\right).
\end{equation}
Because the projection map is idempotent, $P(P(A))=P(A)$, the constraint is equivalent to $P(R)=R$. Consequently, every solution is obtained by choosing $A=R$, although any seed $A=R+N$ with $P(N)=0$ yields the same PO-source.

The projection also satisfies $P(A)^{\ddagger}=P(A^{\ddagger})$, ensuring that all solutions obeying the reality condition are obtained from $A=(Z+Z^{\ddagger})/2$ with arbitrary $Z\in M_8(\mathbb{C})$. Gauge covariance is preserved by covariant seeds with the represented Clifford frame transformed consistently; in a fixed frame, the allowed transformations must preserve $\hat{U}$ and $K$.

The orthogonal complement of a timelike propagation vector is five-dimensional, giving $\binom{5}{r}$ independent grade-$r$ components for $0\leq r\leq5$, while the grade-six component vanishes, leaving a kernel of $32$ complex dimensions in the complexified algebra before imposing reality, divided equally between even and odd parity. As detailed in Appendix~\ref{appx_spinor_lagrangian}, the grade-zero, grade-one, and grade-two solutions within the Plane-of-Now algebra generated by $\{\mathbf{e}_{5},\mathbf{e}_{6}\}$ give rise naturally to mass, isospin, and scalar Yukawa interactions, respectively.

\subsection{Anticommuting Odd-Grade Solutions}
For an odd seed $A_{\mathrm{o}}$ satisfying $KA_{\mathrm{o}}K=-A_{\mathrm{o}}$, the general solution is
\begin{equation}\label{eq_odd_matrix_source}
    R_{\mathrm{o}}=\frac{1}{2}f(\tau)\left(A_{\mathrm{o}}-\hat{U}A_{\mathrm{o}}\hat{U}\right),\quad \{U,R_{\mathrm{o}}\}=0,
\end{equation}
where $f(\tau) \in \mathbb{R}$ is an arbitrary real amplitude that can depend on the spacetime coordinates.

This sector contains the grade-one, grade-three, and grade-five source components, with respective complex dimensions $5$, $10$, and $1$ before imposing reality, and its matrices exchange the propagation eigenspaces associated with the eigenvalues $+c$ and $-c$.

\subsection{Commuting Even-Grade Solutions}
For an even seed $A_{\mathrm{e}}$ satisfying $KA_{\mathrm{e}}K=A_{\mathrm{e}}$, the general solution is
\begin{equation}\label{eq_even_matrix_source}
    R_{\mathrm{e}}=\frac{1}{2}f(\tau)\left(A_{\mathrm{e}}+\hat{U}A_{\mathrm{e}}\hat{U}\right),\quad [U,R_{\mathrm{e}}]=0.
\end{equation}
This sector contains the grade-zero, grade-two, and grade-four source components, with respective complex dimensions $1$, $10$, and $5$, and preserves each propagation eigenspace.

Every solution decomposes uniquely as $R=R_{\mathrm{e}}+R_{\mathrm{o}}$, where
\begin{equation}
    [U,R]=2UR_{\mathrm{o}},\quad \{U,R\}=2UR_{\mathrm{e}}.
\end{equation}
Thus, the general noncommuting class consists of these sums with $R_{\mathrm{o}}\neq0$, while purely anticommuting solutions additionally require $R_{\mathrm{e}}=0$. Preservation of the Plane-of-Now sectors remains the independent condition $[\mathfrak{i},R]=0$.
